\documentclass{article}
\usepackage[margin=1in]{geometry}

\title{Lab 1B Answers}
\author{Student}
\date{}

\begin{document}
\maketitle

\section*{Part 2 - Training choices (by model)}

\textbf{2-16-3}

I used lr 0.01. It trains but not super good on spirals. Blobs are fine. 

\textbf{2-32-32-3}

Also lr 0.01. This one felt more stable and learned faster than 2-16-3.

\textbf{2-64-64-3}

lr 0.01 worked well. This was one of the best overall for me.

\textbf{2-64-64-64-3}

Still lr 0.01. It can fit hard shapes but training is a bit slower.

\textbf{2-128-128-3}

I changed lr to 0.005 because 0.01 was jumpy. It works well but takes more compute.

\textbf{Quick summary:}
small models are fast but can underfit spirals. very big/deep models can do well but cost more and need better tuning.

\section*{Part 2 - Reflection}

\textbf{1) Why does depth help spirals more than blobs?}

Spirals are very nonlinear, so extra layers help bend the boundary. Blobs are almost linearly separable, so depth is less important.

\textbf{2) Why can wide networks work with low depth?}

A wide layer has many neurons, so it can still build complex boundaries even without many layers.

\textbf{3) Which dataset was hardest?}

Spirals. The classes are mixed around each other, so it is harder to separate.

\textbf{4) Effect of depth/width on training?}

More depth/width usually helps fitting, but training can be slower and more sensitive.

\textbf{5) Do deeper networks always do better?}

No. If data is easy, deeper is not always better. Sometimes it just adds complexity.

\section*{Part 3 - Reflection}

\textbf{1) Why no nonlinearity means linear model?}

Because stacking linear layers is still just a linear function.

\textbf{2) What if ReLU backward is wrong for negative values?}

Then gradients are wrong and training can go in bad directions.

\textbf{3) Why use \texttt{torch.autograd.Function}?}

It is easier and less error-prone. PyTorch can connect your custom op to the rest of the graph automatically.

\end{document}
